\subsection{Uniform positivity and derivative estimates}

We now prove the remaining assertions of Theorem~\ref{density:main}.
The main point is to choose one derivative constant \(K_*\), followed
by one threshold for \(Q\), that works at every stage.  A phase
derivative of the correction in \eqref{density:recursion} involves
\(S_p''\).  We obtain a bound for this derivative from the frequency
support and the estimates already proved at earlier stages.

\begin{lemma}[Growth along phase orbits]\label{density:orbit-growth}
Suppose \(Q\ge2\), \(k_j=Q^j\), and the positive functions
\(W_0=1,W_1,\ldots,W_n\), with \(n\ge1\), satisfy
\[
 c_0W_{j-1}\le W_j\le2W_{j-1},
 \qquad |TW_j|\le K_*k_jW_j
 \quad(1\le j\le n).
\]
Put
\[
 R=\frac2{c_0},\qquad \beta=\log_2R.
\]
For every \(p\in S\), the function
\(S_p(\theta)=W_n(e^{i\theta}p)\) satisfies
\begin{equation}\label{density:orbit-growth-bound}
 S_p(\theta+u)
 \le e^{K_*}R(1+k_n|u|)^\beta S_p(\theta),
 \qquad |u|\le\pi.
\end{equation}
In particular, the constants in this bound are independent of \(Q,n,p\).
\end{lemma}

\begin{proof}
Choose the largest \(1\le l\le n\) such that \(k_l|u|\le1\),
or put \(l=0\) if no such index exists.  The comparisons between
successive densities give
\[
 \frac{W_n(e^{i(\theta+u)}p)}{W_n(e^{i\theta}p)}
 \le R^{n-l}
 \frac{W_l(e^{i(\theta+u)}p)}{W_l(e^{i\theta}p)}.
\]
When \(l\ge1\), logarithmic differentiation and the phase derivative
bound show that the last ratio is at most
\(\exp(K_*k_l|u|)\le e^{K_*}\).
For \(l=0\) it equals one.
If \(l=n\), this already proves the assertion.
Otherwise the maximality of \(l\), including the case \(l=0\), gives
\[
 k_n|u|>Q^{n-l-1}\ge2^{n-l-1}.
\]
It follows that
\(R^{n-l}\le R(1+k_n|u|)^\beta\), proving
\eqref{density:orbit-growth-bound}.
\end{proof}

\begin{lemma}[A weighted derivative bound]\label{density:weighted-derivative}
Fix \(A\ge1\) and \(\beta\ge0\).  Let \(q\ge1\) be an integer,
and let \(s\) be a positive trigonometric polynomial of degree at most
\(2q\).  Suppose that
\[
 s(\theta+u)\le A(1+q|u|)^\beta s(\theta)
 \qquad(\theta\in\mathbb R,\ |u|\le\pi).
\]
There is a constant \(C(A,\beta)\), independent of \(q,s\), such that
\[
 |s''(\theta)|\le C(A,\beta)q^2s(\theta).
\]
\end{lemma}

\begin{proof}
Choose a fixed smooth compactly supported function
\(\chi:\mathbb R\to\mathbb R\) which equals one on \([-2,2]\), and put
\[
 K_q(u)=\sum_{h\in\mathbb Z}\chi(h/q)e^{ihu}.
\]
For normalized convolution on the circle, the frequency assumption
implies \(s''=s*K_q''\).  We claim that, for every positive integer \(N\),
\begin{equation}\label{density:kernel-decay}
 |K_q''(u)|\le C_Nq^3(1+q|u|)^{-N},
 \qquad |u|\le\pi,
\end{equation}
where \(C_N\) depends only on \(N,\chi\).

Indeed the coefficients of \(K_q''\) are
\[
 a_h=-h^2\chi(h/q)=q^2g(h/q),
 \qquad g(v)=-v^2\chi(v).
\]
There are \(O(q)\) nonzero coefficients, each of size \(O(q^2)\),
so direct summation gives \(|K_q''|\le Cq^3\).
An \(N\)-th finite difference of \(a_h\) is an iterated integral of
the \(N\)-th derivative of \(q^2g(v/q)\), and hence has size at most
\(C_Nq^{2-N}\).  Only \(O_N(q)\) such differences can be nonzero,
since \(g\) is compactly supported and \(q\ge1\).  Their absolute sum
is therefore at most \(C_Nq^{3-N}\).
Multiplication of the Fourier series by \((1-e^{iu})^N\) produces
these differences, up to an index shift.  Using
\(|1-e^{iu}|\ge2|u|/\pi\) on \([-\pi,\pi]\) gives the second estimate
\[
 |K_q''(u)|\le C_Nq^{3-N}|u|^{-N}\qquad(u\ne0).
\]
Combining it with the direct estimate proves
\eqref{density:kernel-decay}.

Choose an integer \(N>\beta+1\).
Positivity, the assumed growth bound, and
\eqref{density:kernel-decay} now give
\[
 \begin{split}
 |s''(\theta)|
 &\le\int_{-\pi}^{\pi}s(\theta-u)|K_q''(u)|
                         \,\frac{du}{2\pi}\\
 &\le AC_Nq^3s(\theta)
       \int_{-\pi}^{\pi}(1+q|u|)^{\beta-N}\,\frac{du}{2\pi}\\
 &\le C(A,\beta)q^2s(\theta).
 \end{split}
\]
\end{proof}

Combining these two lemmas with the bandwidth assertion of
Lemma~\ref{density:moments} gives the following consequence.
If the estimates of Theorem~\ref{density:main} hold through stage
\(j-1\), then for \(k=k_j\) and the central function \(S_p\) in
\eqref{density:recursion},
\begin{equation}\label{density:central-derivatives}
 \frac{|S_p'|}{kS_p}\le\frac{K_*}{Q},
 \qquad
 \frac{|S_p''|}{k^2S_p}\le\frac{C_1}{Q^2},
\end{equation}
where \(C_1\) depends only on \(K_*,c_0\) and the fixed cutoff
\(\chi\).  It is independent of \(Q,j,p\).
For \(j\ge2\), apply the lemmas with \(q=k_{j-1}\), noting that
\(\ell_{j-1}\le2k_{j-1}\); the first inequality follows directly
from the phase derivative bound.  For \(j=1\), both derivatives vanish
because \(S_p=1\).

We next estimate a single correction using only normalized bounds.
This will allow us to choose \(K_*\) before requiring
\eqref{density:central-derivatives} to be small.
All horizontal vectors \(Z\) below have Euclidean length one.
Enlarge a constant \(D_R\ge1\), depending only on \(R_0\), so that
on every patch at every sufficiently large degree \(k\),
\begin{equation}\label{density:patch-directions}
 |Yy|\le D_R,\qquad |Y\theta|\le D_R/k,\qquad
 |Y(e^{ik\theta})|\le D_R,
 \quad Y=Z/\sqrt{k}\ \hbox{or}\ Y=T/k.
\end{equation}
Here is a direct verification.  With \(c=\langle\xi,p\rangle\),
complex horizontality gives
\[
 |Zc|\le\sqrt{1-|c|^2}
 \le\sqrt{2R_0/k},\qquad |c|\ge e^{-R_0/k}.
\]
Differentiate \(y=-k\log|c|\) and \(\theta=\arg c\) to obtain
the first two bounds for \(Z/\sqrt{k}\); the third follows from
the second.  For \(T/k\), one has \(Yy=0\) and \(Y\theta=1/k\).
The fixed lower threshold for \(k\) makes \(e^{-R_0/k}\ge1/2\).
We also choose \(D_R\) large enough to bound the length of the radial
arc from \(e^{i\theta}p\) to each patch point by \(D_R/\sqrt{k}\).
In adapted coordinates this arc is
\[
 u\longmapsto(e^{i\theta}\cos u,e^{i\gamma}\sin u),
 \qquad 0\le u\le \arccos(e^{-y/k}).
\]
It has unit speed and is complex-horizontal.  The length bound follows
from \(0\le y\le R_0\) and \(1-e^{-y/k}\le R_0/k\).
Fix this enlarged \(D_R\) throughout the rest of the proof.

\begin{lemma}[Normalized correction estimate]
\label{density:normalized-patch}
Write
\[
 F=\|f\|_\infty,\quad B=\|b\|_\infty,\quad
 F_1=\|f'\|_\infty,\quad B_1=\|b'\|_\infty.
\]
On a degree-\(k\) patch suppose \(W_{j-1}>0\), \(S_p>0\), and
\[
 \frac{e^{2y/k}S_p}{W_{j-1}}\le2,\qquad
 r_1=\frac{|S_p'|}{kS_p}\le1,\qquad
 r_2=\frac{|S_p''|}{k^2S_p}\le1.
\]
For the correction \(A_{j,p}\) in \eqref{density:recursion},
\[
 |YA_{j,p}|\le C_2W_{j-1},
 \quad Y=Z/\sqrt{k}\ \hbox{or}\ T/k,
 \qquad
 C_2=2D_R\bigl[4(F+B)+F_1+B_1\bigr].
\]
In particular, \(C_2\) is independent of \(K_*,Q,j,p\).
\end{lemma}

\begin{proof}
Besides \eqref{density:patch-directions}, use
\[
 |Y(e^{2y/k})|\le(2D_R/k)e^{2y/k},\qquad
 |YS_p|\le(D_R/k)|S_p'|,\qquad
 |YS_p'|\le(D_R/k)|S_p''|.
\]
Differentiating \eqref{density:recursion} and taking absolute values gives
\[
 |YA_{j,p}|
 \le D_Re^{2y/k}S_p
 \left[(1+2/k)(F+Br_1)+F_1+B_1r_1+Fr_1+Br_2\right].
\]
Since \(k\ge1\), the assumed normalized bounds imply the stated estimate.
\end{proof}

\begin{proof}[Completion of the proof of Theorem~\ref{density:main}]
Fix \(K_*>2C_2/c_0\).  The constant \(C_1\) in
\eqref{density:central-derivatives} is now fixed as well.
We give a finite list of requirements on \(Q\), all independent of
the induction stage.

The horizontal derivative estimate at stage \(j-1\), integrated along
the radial arc just described, will imply
\begin{equation}\label{density:central-comparison}
 e^{-K_*D_R/\sqrt Q}
 \le\frac{S_p(\theta)}{W_{j-1}(\xi)}
 \le e^{K_*D_R/\sqrt Q}.
\end{equation}
Indeed the logarithmic derivative along a unit horizontal arc is at most
\(K_*\sqrt{k_{j-1}}\), and the arc length is at most \(D_R/\sqrt{k_j}\).
At the initial stage the ratio is exactly one.
Since \(k_j\ge Q\), put
\[
 H(Q)=\exp\left(K_*D_R/\sqrt Q+2R_0/Q\right).
\]

Choose a threshold \(Q_0\) such that every integer \(Q\ge Q_0\)
satisfies all earlier degree and patch requirements, and also
\begin{equation}\label{density:Q-conditions}
 \begin{gathered}
 Q\ge2,\qquad H(Q)\le2,\qquad
 K_*/Q\le1,\qquad C_1/Q^2\le1,\qquad
 Q^{-1/2}\le c_0/2,\\
 H(Q)(F+BK_*/Q)\le1-2c_0.
 \end{gathered}
\end{equation}
These conditions are compatible: \(H(Q)\to1\) and
\(F=1-4c_0<1-2c_0\).

We now fix any integer \(Q\ge Q_0\) and induct on \(j\).
Assume the claimed estimates through stage \(j-1\); at \(j=1\)
the preceding density is the constant one.
Equation~\eqref{density:central-derivatives}, obtained from earlier
stages only, and the horizontal comparison
\eqref{density:central-comparison} establish the normalized hypotheses
of Lemma~\ref{density:normalized-patch}.  Moreover,
\[
 |A_{j,p}|
 \le e^{2y/k}S_p(F+Br_1)
 \le H(Q)(F+BK_*/Q)W_{j-1}
 \le(1-2c_0)W_{j-1}.
\]
At any point at most one patch correction is nonzero.  Hence
\[
 2c_0W_{j-1}\le W_j\le(2-2c_0)W_{j-1},
\]
which proves positivity and the required multiplicative comparison.

For \(Y=Z/\sqrt{k_j}\), the previous derivative bound gives
\(|YW_{j-1}|\le(K_*/\sqrt Q)W_{j-1}\).
For \(Y=T/k_j\), it gives the smaller coefficient \(K_*/Q\).
Together with Lemma~\ref{density:normalized-patch},
\eqref{density:Q-conditions}, and the choice of \(K_*\), these yield
\[
 |YW_j|\le(C_2+K_*/\sqrt Q)W_{j-1}
 \le K_*c_0W_{j-1}\le K_*W_j.
\]
Outside the patches the correction vanishes.  Its smooth zero
extension at their boundaries gives the same estimate there.
Multiplying by \(\sqrt{k_j}\) or \(k_j\) proves the two derivative
bounds and completes the induction.

Lemma~\ref{density:moments} supplies the exact moments, rotation mean
and bandwidth independently of these inequalities.  In particular
\(\int_SW_j\,d\sigma=1\); positivity therefore makes
\(W_j\,d\sigma\) a probability measure.
\end{proof}

We emphasize the order of the choices.  The normalized correction
constant \(C_2\) depends only on the fixed patch data.  It determines
\(K_*\); this in turn determines \(C_1\); finally \(Q_0\) is chosen.
No lower bound for \(W_j\) uniform in \(j\) is required.  The positive
comparison at each stage suffices for logarithmic differentiation,
and all later integral estimates will use the probability normalization.
